\chapter{Conclusions and Future Work}
\label{ch:conclusion}

This thesis investigated whether a tiered edge architecture can divide
visual hazard-detection responsibilities between two heterogeneous,
resource-constrained drone platforms. The Sentinel was assigned initial hazard
detection, temporal filtering and target designation; the Scout was assigned
secondary visual hazard verification. The work combined dataset curation,
six-model screening, deployment conversion, device-specific integration,
ROS~2/DDS communication, stationary two-node trials and development
simulation. The evidence supports a functioning division of responsibilities,
but only within the completed offline, stationary and simulated conditions.

\section{Contributions and principal findings}
\label{sec:conclusion:1}

The first contribution is an explicit tiered edge architecture whose
roles, messages and decision states are traceable. A Sentinel alert carries a
unique event identity and image-space context; a Scout verdict returns the
corresponding identity, outcome and score; and the Sentinel records the final
decision. This interface converts two detector executions into a correlated
workflow and makes confirmation, rejection, timeout and degraded behaviour
observable. The contribution is architectural and implementation-oriented: it
does not by itself prove that a two-stage system is more accurate or reliable
than a single detector.

The second contribution is an evidence-controlled model-development process.
The curated fire-only study froze source-family-aware splits, screened six
configurations, selected operating points from validation data and reserved the
test split for one-time evaluation. YOLOv8n-416 led the validation comparison
with AP50--95 of 0.236 and F1 of 0.612, while EfficientDet led the one-time test
comparison with AP50--95 of 0.168 and F1 of 0.442. The reversal, overlapping
grouped-bootstrap intervals and single training seed preclude a universal
ranking. Conversion fidelity added another selection constraint: both
416-pixel YOLO exports passed the frozen match-rate gate, whereas the two
640-pixel YOLO exports and the SSD export did not. Predictive ranking, runtime
availability and deployment fidelity must therefore be evaluated separately.

The third contribution is the integration and profiling of the historical
deployed pair. The final Sentinel TensorRT path used a 640-pixel engine and
averaged 58.151~ms in the retained camera-free 100-iteration benchmark. The
Scout TFLite path averaged 999.431~ms under its corresponding protocol. These
measurements establish executable platform-specific paths, not a controlled
model or hardware comparison. The older 416-pixel Sentinel benchmark remains a
separate historical result and is not attributed to the final deployment.
Decoder profiling further showed that software around the neural network can
dominate end-to-end cost: vectorization reduced the historical Sentinel
complete-pipeline mean from 204.324 to 57.049~ms, while the Scout remained
inference-bound.

This deployed decision path was onboard in the computational sense: each
vehicle acquired and processed its own camera stream, and the Sentinel hosted
coordination and final-decision handling. The laboratory laptop supported
configuration, observation and evidence retrieval but was not required for
inference. This property strengthens the edge-computing contribution without
changing the experimental boundary: the accepted physical trials remained
stationary, and aerial mission execution remained simulated.

The fourth contribution is a bounded stationary demonstration of secondary
verification. The Scout confirmed 13 of 13 correlated requests in a marked
stronger-flame interval, confirmed one of ten in a smaller unstable-flame
interval and rejected six of six forced unlit requests. Mean
alert-to-final-decision time was approximately 1.44--1.46~s in the accepted
positive conditions, with Scout inference accounting for most of the delay.
These are condition-specific observations, not estimates of industrial
accuracy. The forced unlit requests bypassed Sentinel detection, and the logs
did not retain the complete pre-filter candidate sequence needed to measure the
effect of the temporal filter.

The fifth contribution is a documented coordination and simulation boundary.
Two 300-event ROS~2/DDS campaigns returned all application-level replies, with
median round-trip times of 5.878 and 5.257~ms; the earlier campaign also
retained a 400.453~ms tail. A separate 1,000-event campaign under concurrent
live inference load delivered every reply with 7.984~ms p95 RTT. Four
five-repetition hardware fault cases demonstrated fail-closed camera-interface,
stale-stream and unavailable-Scout paths plus idempotent rejection of ten
post-timeout verdicts. Development simulation completed two nominal and four
selected-fault pilots and exposed a mission-state defect that was corrected and
rerun. One later hash-linked visual run completed the confirmed-hazard sequence
from patrol through Scout dispatch, verification, alarm publication and landing
with distinct Sentinel and Scout proxies. Its final continuous delivery was
reviewed and accepted for thesis and defence presentation, while the synchronized
logs remained authoritative for event timing. These bounded cases do not provide
a general reliability estimate, and the incomplete formal campaign does not
support a claim of physical mission correctness.

\section{Assessment of the central hypothesis}
\label{sec:conclusion:2}

The central hypothesis is supported to a limited and conditional degree. The
tested heterogeneous edge nodes can execute different perception roles, exchange
correlated events and produce both confirmations and rejections in a stationary
two-node workflow. The results also demonstrate that the division has a
measurable cost: secondary verification adds roughly a second of Scout
inference and increases final-decision latency. Typical communication time was
small relative to that inference cost under the retained laboratory network,
although the observed tail prevents a general low-latency guarantee.

The evidence does not establish four stronger interpretations of the
hypothesis. First, the complete stationary service configuration survived a
registered 30-minute run with both perception paths, communication, logging and
flight services active, but it did not satisfy every criterion: the Scout p95
alert-to-verdict time was 2.618~s and Sentinel swap increased by 0.922~MiB.
This supports bounded coexistence, not unlimited or airborne sustainability.
Second, no
common-sequence ablation measured whether the temporal filter and Scout stage
reduce false alerts while retaining labelled hazards. Third, the simulation
pilots and stationary physical trials do not demonstrate autonomous
two-vehicle flight, geographic localization or industrial-site operation.
Fourth, transfer performance was not evaluated because neither external data
source reached a frozen eligible inference manifest.

The answer to the research problem is consequently narrower than the intended
application. The project establishes a traceable prototype and a reproducible
evaluation framework for heterogeneous role division. It does not establish a
validated industrial hazard-response system. In particular, the final physical
evidence concerns visible flame appearance in RGB imagery; physical smoke,
heat, gas and particulate sensing were not evaluated.

\section{Prioritized future work}
\label{sec:conclusion:3}

Future work should first close the evidence gaps that prevent answers to the
existing research questions. Architectural extensions should follow only after
the current prototype has passed those checks.

\subsection{Complete the registered validation programme}

The first priority is the strengthened candidate-by-device screen. Each model
admitted to a device must first pass the frozen conversion-fidelity gate, then
complete the registered warm-up and measured iterations with artifact,
configuration and runtime hashes. The completed 30-minute full-service study
should then be repeated across controlled ambient and power conditions, adding
calibrated power measurement and a sustained tiered-workflow load. One-second
resource records should cover memory, temperature, CPU/GPU load, power, dropped
frames, queue depth and service failures. These tests are required to answer
RQ1 fully and to strengthen the conditional answer to RQ2.

The second priority is a common-stimulus stage ablation. A frozen labelled
sequence should be replayed through single-frame Sentinel detection, the
three-consecutive-frame temporal filter and Scout verification. The analysis
should report event retention, false-alert burden, false-negative burden and
decision latency at every stage, with the experimental unit and uncertainty
method defined in advance. This directly answers the unresolved part of RQ3
without inferring filter performance from incomplete live logs.

The third priority is to complete the registered simulation campaign using
frozen airframe parameters, scenario seeds, mission timeouts and pass/fail
rules. Failures and retries must remain in the retained denominator. After the
software and simulation gates pass, physical work should proceed in stages:
bench integration with both flight controllers, restrained or propeller-free
tests, controlled single-vehicle flight and finally supervised two-vehicle
missions under approved safety procedures. This sequence is necessary before
RQ4 can extend beyond development simulation.

The fourth priority is external evaluation. FASDD-UAV requires acquisition,
licence and provenance checks. MmodalFire requires closure of timestamp,
event-label, overlap and eligibility gates. Only a frozen eligible manifest
should enter inference. Results should be reported per source domain and target
scale, with overlap exclusions and unavailable conditions visible, so that RQ5
measures transfer rather than hidden data relatedness.

\subsection{Improve perception, localization and operational readiness}

After the validation programme, perception should expand beyond a second RGB
fire detector. Calibrated thermal, particulate or gas sensing would supply a
measurement with a different physical basis and reduce dependence on correlated
visual errors. New data collection should cover smoke, larger standoff
distances, small targets, glare, clouds, occlusion, motion blur and industrial
backgrounds. Repeated frozen seeds and independent evaluation sources would
separate model-selection effects from training and split variability.

The current target designation must also be converted from image-space metadata
to a shared spatial estimate. This requires camera calibration, vehicle pose,
time synchronization, frame transforms, uncertainty propagation and a validated
safe-standoff planner. The Scout should use the resulting region or world
estimate for directed observation rather than full-frame inference alone.

Operational readiness additionally requires hazard-specific alarm deadlines,
human oversight, fail-safe policies, persistent artifact staging, radio
planning, data governance, environmental qualification and compliance with
applicable aviation and site procedures. These are requirements for future
evaluation, not properties of the current implementation.

\reviewstatus{No new result is introduced. RQ2 has a bounded answer and RQ5 was not evaluated;
physical smoke verification, geographic localization, autonomous industrial
flight and universal reliability remain outside the evidence.}
